\chapter{European and translational alignment}
\label{chap:european_and_translational_alignment}

European-facing terminology requires exact anatomical and pathology evidence before using labels such as oral cavity squamous cell carcinoma, HNSCC or HPV-related disease. These fields are unavailable, so the broad oral-cancer label is retained. Matched non-tumour is not called healthy or normal.

For future translation, \textcite{eu2017ivdr} distinguishes scientific validity, analytical performance and clinical performance for an intended purpose. This discovery matrix contributes exploratory scientific evidence only. A future assay would require locked analytes, measurement procedures, detection limits, precision, interference assessment, reference standards where possible, intended-use participants and independent clinical performance evidence. Phase 7's protocol is preparatory, not a conformity assessment.

BRISQ supports complete reporting of biospecimen collection, processing and storage \parencite{moore2011}. MIAPE supports proteomics acquisition and processing metadata \parencite{taylor2007}. Both reveal gaps rather than certify the present data. Ethics approval, consent scope, lawful basis, data-controller roles and permitted secondary use must be recovered before sharing; pseudonymised health and molecular data should not be assumed anonymous.

TRIPOD and PROBAST shape prediction reporting and risk-of-bias interpretation \parencite{collins2015,wolff2019}. Under PROBAST reasoning, the single small cohort, discovery-only participants, incomplete clinical metadata and absent external validation create high applicability concern despite strong internal discrimination.
